\section{Method}

\subsection{Task Formulation}
\label{sec:task-formulation}

We frame our task as a Partially Observable Markov Decision Process (POMDP). Each rollout trajectory is divided into a sequence of trajectory units indexed by $i$. A trajectory unit is defined as 
$U_i = (G_i, O_i, S_i, A_i, A_p, O_p)$, 
where $G_i$ denotes the task instruction, and $O_i$, $S_i$, and $A_i$ denote the observation, robot state, and action chunks, respectively, accompanied by an action placeholder $A_p$ and an observation placeholder $O_p$.

Each chunk (i.e., $O_i$, $S_i$, and $A_i$) consists of a sequence of time-indexed elements, including observations, states, or actions extracted from the raw trajectory $\tau$. Each element (i.e., $o_t$, $s_t$, $a_t$) is further tokenized into a sequence of tokens that serves as the actual model input.

In unit~$i$, the model first predicts an action chunk $\hat{A}_i$
conditioned on the instruction $G_i$, the observation chunk $O_i$, 
and the state chunk $S_i$.
It then predicts the next observation chunk conditioned on 
$(O_i, S_i, A_i)$ during training or $(O_i, S_i, \hat{A}_i)$ during inference\footnote{A history buffer $h$ is also included in the conditioning; detailed formulation is provided in a later section.}.
Executing the predicted actions advances the environment and produces 
the next unit $U_{i+1}$.
A summary of the key notations used in the WorldAgen framework is 
provided in Table~\ref{tab:notation} in the Appendix.
\subsection{Joint State--Action Prediction}
\label{sec:join-action}

WorldAgen contains two predictive components: a task-conditioned agent model and a task-agnostic world model. Both models share a single Transformer backbone, enabling a unified representation across tasks, as illustrated in Figure~\ref{fig:worldagen_overview}(a).

We maintain two types of histories:

\textbf{Agent history:}
\[
h^a_i = (G_i, O_x, S_x, \hat{A}_x)_{x=i-k}^{i-1},
\]
which contains the task instruction $G_i$ along with the observation, robot state, and predicted action chunks from the previous $k$ segments. This history is used by the agent model for task-aware action prediction.

\textbf{World history:}
\[
h^w_i = (O_x, S_x, \hat{A}_x)_{x=i-k}^{i-1},
\]
which includes the observation, robot state, and predicted action chunks from the previous $k$ segments, and is used by the world model to predict task-agnostic environment dynamics.

We define the two models at trajectory segment $U_i$ as follows:

\textbf{Agent Model.}
The agent model predicts the next action chunk $\hat{A}_i$ conditioned on the current observation chunk $O_i$, robot state chunk $S_i$, task instruction $G_i$, and agent history $h^a_i$:
\[
p_a\!\left(
  \hat{A}_i
  \,\big\vert\,
  G_i, O_i, S_i, h^a_i
\right).
\]

\textbf{World Model.}
The world model predicts the next observation chunk $\hat{O}_i$ conditioned on the current observation chunk $O_i$, robot state chunk $S_i$, predicted action chunk $\hat{A}_i$, and world history $h^w_i$:
\[
p_w\!\left(
  \hat{O}_i
  \,\big\vert\,
  O_i, S_i, \hat{A}_i, h^w_i
\right).
\]

During training, we adopt teacher forcing and replace all predicted actions $\hat{A}_i$ with the ground-truth actions $A_i$.


\begin{figure*}[t]
    \centering
    \includegraphics[scale=0.55]{figures/figure2.pdf}
    \caption{
\textbf{Trajectory Splitting, Chunking, and Two-Step Inference.}
Example with observation chunk length $n=2$ and action chunk length $m=3$.
(a) \textit{Trajectory Splitting and Chunking:}
a raw trajectory is divided into trajectory units $U_i$, where each unit
contains an observation chunk $O_i$ and a robot state chunk $S_i$ obtained
via uniform subsampling from $m$ consecutive observations and states, and an
action chunk $A_i$ obtained by grouping $m$ actions.
(b) \textit{Two-Step Inference:}
the agent model first initializes a zero action placeholder $A_p$ and then
replaces it with the predicted action chunk $\hat{A}_i$.
The world model subsequently predicts the next observation chunk
$\hat{O}_i$, conditioned on the updated trajectory unit.
}
    \label{fig:worldagen}
\end{figure*}



\subsection{Trajectory Splitting and Action Chunking}
\label{traj_split}
Predicting multiple actions at once has been shown to
improve the execution efficiency of real-world robots \cite{intelligence2025pi_}. However, consecutive
observation frames often contain highly redundant visual information, and predicting long observation sequences slows down both training and deployment. To balance these factors, WorldAgen supports different action chunk lengths $m$ and observations chunk length $n$.
We further propose a robust trajectory splitting and action chunking scheme that flexibly adapts to different experimental configurations.

Concretely, given a rollout $\tau = \{(o_t, s_t, a_t)\}_{t=0}^{T-1}$ with $T$ timesteps, we
transform it into a sequence of trajectory units $\{U_i\}_{i=0}^{I-1}$ with $I$ units.
The initial unit $U_0$ aggregates the first $n$ timesteps of the rollout and is paired
with the first action chunk $A_0 = A_{0 \rightarrow n-1} = [a_0, a_1, \ldots, a_{n-1}]$.
For units $U_i$ with $i>0$, we uniformly sample $n$ timesteps to construct the current
observation chunk $O_i$, state chunk $S_i$, and associated action chunk $A_i$:
\begin{align*}
O_i &= O^{n}_{n+(i-1)m \rightarrow n+im-1} \\
    &= \big[
             o_{\,n+(i-1)m},\,
             o_{\,n+(i-1)m + \lfloor \tfrac{1}{n} m \rfloor},\,
             \ldots,\,
             o_{\,n+(i-1)m + \lfloor \tfrac{n-1}{n} m \rfloor}
       \big], \\
S_i &= S^{n}_{n+(i-1)m \rightarrow n+im-1} \\
    &= \big[
          s_{\,n+(i-1)m},\,
          s_{\,n+(i-1)m + \big\lfloor \tfrac{1}{n} m \big\rfloor},\,
\\
    &\quad
          \ldots,\,
          s_{\,n+(i-1)m + \big\lfloor \tfrac{n-1}{n} m \big\rfloor}
       \big], \\
A_i &= A_{n+(i-1)m \rightarrow n+im-1} \\
    &= \big[
             a_{\,n+(i-1)m},\,
             a_{\,n+(i-1)m+1},\,
             \ldots,\,
             a_{\,n+im-1}
       \big].
\end{align*}
This procedure keeps each observation chunk fixed to length $n$, reducing redundant
visual information and improving runtime efficiency.


Figure~\ref{fig:worldagen} (a) illustrates an example with $n=2$ and $m=3$. Algorithm of trajectory splitting and action chunking for variable length of action, observation and robot state chunk is shown in Algorithm~\ref{alg:traj-splitting} in the Appendix, where $\textsc{UniformSubsample}(a,b,n)$ uniformly samples $n$ distinct indices from the integer interval $\{a,\ldots,b\}$ without replacement and returns them in ascending order, which we then use to construct the
observation and state chunks.


\begin{figure}[t]
    \centering
    \includegraphics[scale=0.5]{figures/attn_mask.pdf}
    \caption{
Proposed Mixed Unidirectional Attention Mask: combines local and global masking. For each trajectory unit, WorldAgen applies a local mask that prevents the action chunk $A_p$ from attending to the corresponding action placeholder $A_i$. A global mask is further used to construct a task-agnostic world model head, so that the observation placeholder $O_p$ will not attend to the task instruction $G_i$. General structure conforms to strictly causal prediction.
}
    \label{fig:attn_mask}
\end{figure}

\subsection{Mixed Unidirectional Attention Mask}
To ensure that the world modeling and agent policy tasks are trained jointly without information leakage, we introduce a mixed unidirectional attention masking mechanism as shown in Figure \ref{fig:attn_mask}. 
This mechanism integrates two complementary components:

\textbf{Local mask:} ensures that the current observation chunk $O_i$, the robot proprioception chunk $S_i$, the task instruction $G_i$, and the action placeholder $A_p$ cannot access the action chunk $A_i$. This prevents information leakage within a single timestep.

\textbf{Global mask:} guarantees that the world model head remains invisible to the task instruction $G_i$ across all time steps. Additionally, the global mask enforces mutual invisibility among the task instruction and observations to avoid cross-modal leakage. 

Meanwhile, the general structure conforms to strictly causal prediction, such that both the world model and policy heads cannot access any future tokens.

\subsection{Pretraining}
We pretrain WorldAgen using large-scale trajectory data collected from robot demonstrations. We split trajectories using our proposed robust splitting method.
During pretraining, we jointly optimize the agent model and world model using teacher forcing:

\[
\mathcal{L} = \mathcal{L}_a + \lambda \mathcal{L}_o,
\]
where $\mathcal{L}_a$ is the cross-entropy loss for predicting future action chunks, $\mathcal{L}_o$ is the reconstruction loss for predicting future observation chunks, and $\lambda$ is a weighting factor balancing the two objectives.

It is worth noting that during training, we feed the ground-truth action chunk (teacher forcing), whereas during testing, the model performs autoregressive prediction by feeding the previously predicted action chunk at each timestep.

\subsection{Two-Step Inference}

At test time, WorldAgen runs in a two-step inference loop over trajectory units as shown in Figure \ref{fig:worldagen} (b). For unit $i$, we first perform the agent model step. We construct
the input sequence with a zero-initialized action placeholder $A_p$ and feed it to the agent model, which predicts the current action chunk $\hat{A}_i$. The predicted actions then overwrite the placeholder $A_p$.

Next, we perform the world model step. Conditioned on the task instruction $G_i$, the current observation chunk
$O_i$ and the updated action chunk $A_i$ (same as $\hat{A}_i$), the world model predicts the next observation chunk $\hat{O}_i$. During testing, these two steps are applied alternately to roll out the whole trajectory.

\subsection{Test-Time Training}
While joint pretraining improves world understanding, distribution shifts in novel environments can still degrade performance.
To address this, we introduce a two-stage Test-Time Training strategy that adapts the world model online, refining its understanding of environment dynamics and indirectly improving downstream action prediction. The whole pipeline is given in Figure~\ref{fig:worldagen_overview} (b).

\paragraph{Stage 1: Random Exploration and Data Collection.}
At deployment, the agent first performs free exploratory rollouts in the target environment.
During this phase, it records trajectories of $(O_i, S_i, A_i)$.
To ensure that adaptation focuses on environment dynamics rather than task-specific instructions, all collected trajectories are relabeled with a generic ``no-lang'' token in place of language instructions.

\paragraph{Stage 2: World Model Adaptation.}
Using the collected trajectories, we adapt only the shared backbone with LoRA-based parameter-efficient fine-tuning, while keeping the policy head frozen.  
The update rule is:
\[
\theta^\prime_w \leftarrow \theta_w - \eta \nabla_{\theta_w} \mathcal{L}_o,
\]
where $\theta_w$ are the parameters of the world model head and $\eta$ is the learning rate.

This targeted update improves world understanding ability and enhances the shared representations, which are in turn leveraged by the agent model for action prediction.

By decoupling adaptation from the task goal and restricting it to the world model, \textbf{our TTT procedure enables robust scene adaptation without requiring additional task-specific annotations, paving the way for scalable test-time adaptation in VLA models.}
